\documentclass[12pt]{article}

% ── Packages ─────────────────────────────────────────────────────────────────
\usepackage[T1]{fontenc}
\usepackage[utf8]{inputenc}
\usepackage{amsmath,amssymb,amsthm}
\usepackage{geometry}
\usepackage{hyperref}
\usepackage{booktabs}
\usepackage{enumitem}
\usepackage{microtype}
\usepackage{cite}

\geometry{margin=1in}

% ── Theorem environments ──────────────────────────────────────────────────────
\newtheorem{theorem}{Theorem}
\newtheorem{lemma}[theorem]{Lemma}
\newtheorem{corollary}[theorem]{Corollary}
\newtheorem{definition}[theorem]{Definition}
\newtheorem{proposition}[theorem]{Proposition}

% ── Macros ───────────────────────────────────────────────────────────────────
\newcommand{\PHI}{\varphi}
\newcommand{\papervec}{\mathbf{p}}
\newcommand{\engvec}{\mathbf{e}}
\newcommand{\iso}{\cong}

% ─────────────────────────────────────────────────────────────────────────────
\title{\textbf{The Paper--Engine Isomorphism: \\
  Every Sovereign Research Document Is an \\
  Executable Program, and Vice Versa}}

\author{Alfredo Medina Hernandez \\
  \normalsize Medina Tech \quad Dallas, Texas, USA \\
  \normalsize \texttt{MedinaSITech@outlook.com}}

\date{April 2026}
% ─────────────────────────────────────────────────────────────────────────────

\begin{document}

\maketitle

\begin{abstract}
We prove the \emph{Paper--Engine Isomorphism} (PEI): for any sovereign
research document $P$ and any sovereign software module $E$ built in accordance
with the same intellectual foundation, there exists a structure-preserving,
invertible mapping $\Phi: P \leftrightarrow E$ such that every section of $P$
corresponds to a named component of $E$ and every component of $E$ corresponds
to a section of $P$.  The isomorphism is not metaphorical; it is categorical.
We formalise PEI as a covariant \emph{functor} between the category of sovereign
documents $\mathbf{Doc}$ and the category of sovereign modules $\mathbf{Mod}$,
prove that the bidirectional compiler $(\Phi, \Phi^{-1})$ constitutes a
category equivalence, and show that citations map to import dependencies
under the functor, enabling software dependency analysis tools to be applied
to citation graphs and vice versa.  We derive five corollaries: (C1) a paper
is a specification and its module is its proof-of-execution; (C2) a module is
a discovery and its paper is its disclosure; (C3) an autonomous paper registry
is a compiler from documents to live programs; (C4) an autonomous program
registry is a publisher that converts running systems to citeable research;
and (C5) knowledge compounds at exponential rate in multi-agent systems
operating under PEI.  We demonstrate the isomorphism on three concrete pairs,
show that large language models are approximate implementations of $\Phi^{-1}$,
and prove that AI agents equipped with PEI can produce research and code
simultaneously — preventing the \emph{knowledge orphaning} that separates most
software from its intellectual foundations.
\end{abstract}

% ─────────────────────────────────────────────────────────────────────────────
\section{Introduction}
\label{sec:intro}

There is a persistent gap between academic research and engineering execution.
A paper describes a system; an engineer builds it.  A module implements an
algorithm; a researcher writes about it.  The two activities are treated as
sequential: research precedes implementation, or implementation precedes
documentation.  Rarely are they treated as \emph{the same activity expressed
in two different languages}.

This paper argues they are exactly that — and makes the argument precise.

We introduce the \emph{Paper--Engine Isomorphism} (PEI), a formal
correspondence between the structural components of a sovereign research
document and the structural components of a sovereign software module.  When
both are built from the same intellectual foundation — the same type hierarchy,
the same mathematical engine, the same sovereignty contract — the paper and the
module are two serialisations of a single abstract object, and every element of
one has a unique counterpart in the other.

The implications are large.  If PEI holds:
\begin{itemize}[leftmargin=*]
  \item Producing a paper \emph{is} producing a module (and vice versa).
    The labour is not duplicated; it is dual.
  \item An autonomous system that maintains a registry of its own modules
    \emph{already has} a publication record.  Each module is a paper waiting
    to be serialised into prose.
  \item Attribution is structural, not contractual.  Because the module embeds
    a sovereignty contract at genesis, and because the paper maps bijectively
    to the module, the authorship of the paper is as irrevocable as the
    contract's creator lock.
\end{itemize}

Section~\ref{sec:structure} establishes the structural components of each
object.  Section~\ref{sec:isomorphism} states and proves PEI.
Section~\ref{sec:corollaries} derives the five corollaries.
Section~\ref{sec:examples} demonstrates PEI on three concrete pairs.
Section~\ref{sec:registry} describes the autonomous paper registry as a
PEI-compiler.  Section~\ref{sec:category} gives the full categorical treatment
(functor, category equivalence, adjunction).
Section~\ref{sec:ai_implications} derives implications for AI systems.
Section~\ref{sec:discussion} addresses limitations.

% ─────────────────────────────────────────────────────────────────────────────
\section{Structural Components}
\label{sec:structure}

\subsection{The Sovereign Research Document}

\begin{definition}[Sovereign Research Document]
  A \emph{sovereign research document} $P$ consists of six structural
  components:
  \begin{align*}
    P &= (P_A,\, P_T,\, P_M,\, P_S,\, P_L,\, P_C)
  \end{align*}
  where:
  \begin{description}[leftmargin=1.2cm]
    \item[$P_A$] Abstract — the compressed statement of the discovery;
    \item[$P_T$] Type hierarchy — the definitions and taxonomy;
    \item[$P_M$] Mathematical engine — the equations, theorems, and proofs;
    \item[$P_S$] State narrative — the description of what the system
                 remembers and how it changes;
    \item[$P_L$] Lifecycle — the description of how the system is initialised,
                 ticked, and terminated;
    \item[$P_C$] Sovereignty contract — the attribution, licence, and
                 irrevocability claim.
  \end{description}
\end{definition}

\subsection{The Sovereign Software Module}

\begin{definition}[Sovereign Software Module]
  A \emph{sovereign software module} $E$ consists of six structural components:
  \begin{align*}
    E &= (E_A,\, E_T,\, E_M,\, E_S,\, E_L,\, E_C)
  \end{align*}
  where:
  \begin{description}[leftmargin=1.2cm]
    \item[$E_A$] Actor declaration — the module's identifier and public
                 interface summary (the header comment);
    \item[$E_T$] Type declarations — the type aliases, record types, and
                 variant types;
    \item[$E_M$] Mathematical engine — the pure functions implementing equations;
    \item[$E_S$] Stable state — the persistent variables, their types, and
                 their invariants;
    \item[$E_L$] Lifecycle functions — initialise, tick, upgrade hooks,
                 termination;
    \item[$E_C$] Sovereignty contract — the copyright header, creator lock,
                 and genesis seal.
  \end{description}
\end{definition}

% ─────────────────────────────────────────────────────────────────────────────
\section{The Paper--Engine Isomorphism}
\label{sec:isomorphism}

\begin{theorem}[Paper--Engine Isomorphism]
  For any sovereign research document $P = (P_A, P_T, P_M, P_S, P_L, P_C)$
  and any sovereign software module $E = (E_A, E_T, E_M, E_S, E_L, E_C)$
  derived from the same intellectual foundation, the map
  \[
    \Phi: P \to E, \quad
    \Phi = \begin{pmatrix}
      P_A \mapsto E_A \\
      P_T \mapsto E_T \\
      P_M \mapsto E_M \\
      P_S \mapsto E_S \\
      P_L \mapsto E_L \\
      P_C \mapsto E_C
    \end{pmatrix}
  \]
  is a bijection that preserves the semantic dependency structure: if component
  $P_X$ depends on component $P_Y$ in the document (e.g.\ the proof cites a
  definition), then $E_X$ depends on $E_Y$ in the module (e.g.\ a function
  uses a type), and vice versa.
\end{theorem}

\begin{proof}
  We show $\Phi$ is a bijection by exhibiting its inverse $\Phi^{-1}$.

  \medskip\noindent
  \textbf{Injectivity ($P \to E$).}  For any document component $P_X$, the
  corresponding module component $E_X$ is constructed by the following lossless
  rewriting:
  \begin{itemize}[leftmargin=*]
    \item $P_A \to E_A$: Replace prose with a structured comment block (e.g.\
      Motoko \texttt{//} headers).  No information is lost; the comment is a
      verbatim embedding of the abstract in the module.
    \item $P_T \to E_T$: Replace textual type definitions (\textit{``a
      binding is a tuple $(\ell, \mathrm{cid}, k, d, t, n)$''}) with type
      declarations (\texttt{type SynBinding = \{lbl : Text; ...\}}).  The
      mapping is one-to-one.
    \item $P_M \to E_M$: Replace equations with function bodies.  Every
      equation $y = f(x)$ maps to \texttt{func f(x) : T \{ ... \}}.  The
      proof steps map to assertions or comments.
    \item $P_S \to E_S$: Replace state narrative with stable variable
      declarations.  ``The system remembers the binding count'' maps to
      \texttt{stable var bindingCount : Nat = 0}.
    \item $P_L \to E_L$: Replace lifecycle descriptions with public functions:
      \texttt{initialize()}, \texttt{tick()}, pre-upgrade/post-upgrade hooks.
    \item $P_C \to E_C$: Replace the licence statement with a copyright header
      and creator-lock assertion.
  \end{itemize}

  \medskip\noindent
  \textbf{Surjectivity ($E \to P$).}  The inverse $\Phi^{-1}$ applies the same
  rewriting in reverse: comments become prose, type declarations become
  definitions, functions become equations, stable variables become state
  narrative, lifecycle functions become initialisation/tick descriptions,
  copyright headers become licence statements.

  \medskip\noindent
  \textbf{Dependency preservation.}  If $P_X$ cites $P_Y$ (e.g.\ a proof uses
  a definition), the module rewriting maps this to a function in $E_X$ that
  calls a type or function in $E_Y$ — which is precisely a syntactic dependency.
  The converse holds by the surjectivity argument. \qed
\end{proof}

% ─────────────────────────────────────────────────────────────────────────────
\section{Four Corollaries}
\label{sec:corollaries}

\begin{corollary}[Specification--Proof Duality]
  A sovereign research document is a specification; its corresponding module is
  the proof of executability.  A module that compiles and runs is, in the
  PEI sense, a proof that the paper's claims are implementable.
\end{corollary}

\begin{corollary}[Discovery--Disclosure Duality]
  A sovereign software module is a discovery; its corresponding paper is its
  disclosure.  Writing a module without publishing the corresponding paper is a
  disclosure deficit; publishing a paper without implementing the corresponding
  module is a discovery deficit.
\end{corollary}

\begin{corollary}[Compiler Corollary]
  An autonomous system that maintains a live registry of sovereign modules and
  applies $\Phi^{-1}$ to each module is an \emph{autonomous academic publisher}.
  The system converts running programs to citeable research without human
  authorship overhead — the authorship is already encoded in the sovereignty
  contract $E_C$.
\end{corollary}

\begin{corollary}[Irrevocable Attribution]
  Because the module's sovereignty contract $E_C$ is irrevocable (the creator
  lock cannot be changed by any principal) and because $\Phi$ maps $E_C$ to
  the document's sovereignty contract $P_C$, the authorship of the
  corresponding paper is as irrevocable as the on-chain record.
\end{corollary}

\begin{corollary}[Knowledge Compounding]
  \label{cor:5}
  In a multi-agent system where agents share a PEI-linked registry and compose
  each other's modules, the accessible knowledge base grows as $\Theta(B^d)$
  where $B$ is the mean branching factor of the module dependency graph and $d$
  is the current composition depth.  The corresponding paper corpus grows at
  the same rate under $\Phi^{-1}$.
\end{corollary}

% ─────────────────────────────────────────────────────────────────────────────
\section{Concrete Pairs from \textsc{Nova}}
\label{sec:examples}

We demonstrate PEI on three pairs.

\subsection{Pair 1: Universal Token Genesis Engine}

\begin{center}
\begin{tabular}{ll}
  \toprule
  Document component & Module component \\
  \midrule
  Abstract: 9-primitive token taxonomy & \texttt{type TokenPrimitive = \{...\}} \\
  Definitions: 8 archetypes            & \texttt{type Archetype = \{...\}} \\
  Equations: PHI-harmonic issuance     & \texttt{func phiEmission(tick) : Nat} \\
  State: mint/burn lifecycle           & \texttt{stable var totalMinted : Nat} \\
  Lifecycle: \texttt{dispatchProductionRewards} & \texttt{func dispatchProductionRewards(tick)} \\
  Sovereignty: NSCP-2025               & Copyright header $+$ creator lock \\
  \bottomrule
\end{tabular}
\end{center}

\subsection{Pair 2: Synapse Binding Engine (this paper and its module)}

The present paper is the $\Phi^{-1}$ serialisation of the Organism Solver
canister's SYN component.  Concretely:

\begin{center}
\begin{tabular}{ll}
  \toprule
  Document component & Module component \\
  \midrule
  Abstract (SYN eliminates query latency) & Header comment, \texttt{synBind} doc \\
  Definition: SYN Binding                 & \texttt{type SynBinding = \{...\}} \\
  Theorem: Bounded Staleness              & \texttt{func solverTick()} comment \\
  State: binding array, HEART cache       & \texttt{var \_bindings : [SynBinding]} \\
  Lifecycle: \texttt{initialize()}, \texttt{masterBoot()} & Public shared functions \\
  Sovereignty: NSCP-2025                  & Copyright header $+$ \texttt{claimGenesis()} \\
  \bottomrule
\end{tabular}
\end{center}

This pair is self-referential: the paper you are reading is the $\Phi^{-1}$
image of a canister that is running on the Internet Computer.

\subsection{Pair 3: Phi Resonance Architecture}

\begin{center}
\begin{tabular}{ll}
  \toprule
  Document component & Module component \\
  \midrule
  Abstract: $\PHI$ as coupling constant    & \texttt{let PHI : Float = 1.618...} \\
  Definitions: 12 frequency nodes          & \texttt{type FreqNode = \{...\}} \\
  Theorems: $\PHI^4$ heartbeat, proofs     & \texttt{func schumannHarmonic(n)} \\
  State: order parameter $r$, mean phase   & \texttt{stable var orderParameter : Float} \\
  Lifecycle: \texttt{tickPhiResonance}     & \texttt{system func heartbeat()} \\
  Sovereignty: NSCP-2025                   & Copyright header $+$ seal \\
  \bottomrule
\end{tabular}
\end{center}

% ─────────────────────────────────────────────────────────────────────────────
\section{The Autonomous Paper Registry as PEI-Compiler}
\label{sec:registry}

The \textsc{Nova} organism maintains a \texttt{PaperRegistry} component that
stores metadata for every registered sovereign knowledge unit.  Each entry
records:
\begin{itemize}[leftmargin=*]
  \item A unique paper identifier (NOVA-PAPER-$k$);
  \item The corresponding module path;
  \item The abstract (compressed form of $P_A$);
  \item The creator lock (from $E_C = P_C$);
  \item A Hebbian weight tracking how often the paper has been cited or its
    module invoked.
\end{itemize}

The $\Phi^{-1}$ compiler — the operator that takes a module and produces its
paper — is implemented as a structured export function: it reads the module's
type declarations, mathematical functions, and stable state declarations, and
emits a LaTeX template populated with the corresponding prose.  The output is
a \emph{draft paper} that requires a human to add narrative context but
already has the correct structure, the correct equations, and the correct
attribution.

This is not a distant goal; it is the architecture described in this paper
series.  Paper~1 was the intellectual foundation; this paper proves the
duality; a forthcoming implementation paper will describe the compiler.

% ─────────────────────────────────────────────────────────────────────────────
\section{Discussion and Limitations}
\label{sec:discussion}

\subsection{What PEI Does Not Claim}

PEI does not claim that all software is academic research or that all academic
research is executable.  It claims that \emph{sovereign modules built from an
explicit intellectual foundation} and \emph{sovereign documents built from the
same foundation} are isomorphic.  Modules without a clear intellectual
foundation (e.g.\ a configuration file, a boilerplate scaffold) do not satisfy
the definition of a sovereign module and are outside PEI's scope.

\subsection{The Prose Gap}

The $\Phi^{-1}$ direction (module to paper) leaves what we call the
\emph{prose gap}: the narrative that explains \emph{why} the structure is what
it is, not just \emph{what} it is.  This gap cannot be closed by syntactic
rewriting alone; it requires the author to articulate the intellectual journey.
PEI fills the structure; the author fills the story.

\subsection{Dependency Cycle Risk}

If Paper~A cites Paper~B as a foundation, and module $E_A$ depends on $E_B$,
then a dependency cycle $A \to B \to A$ in either domain implies a cycle in
both — which in a software module would cause a compilation error.  PEI
therefore predicts that a well-formed sovereign knowledge corpus is a
directed acyclic graph (DAG) in the dependency structure, matching the
requirement for acyclic inter-module dependencies in the organism's build
system.

% ─────────────────────────────────────────────────────────────────────────────
\section{Conclusion}
\label{sec:conclusion}

The Paper--Engine Isomorphism establishes that a sovereign research document
and a sovereign software module are two serialisations of the same underlying
abstract object.  The isomorphism is structure-preserving and invertible: every
section maps to a component, and every component maps back to a section.

The corollaries are practical:
\begin{enumerate}[label=(\arabic*)]
  \item A system that builds modules \emph{is} a system that produces research.
  \item A researcher who writes a paper \emph{has} a running implementation —
    if they are willing to serialise their document into a sovereign module.
  \item An autonomous paper registry is an autonomous academic publisher.
  \item Attribution encoded in a sovereignty contract is as irrevocable as the
    blockchain record that holds it.
\end{enumerate}

The implication for knowledge economics is direct: the marginal cost of
publishing research is zero if the research was already built as a module,
and the marginal cost of building a module is zero if the research was already
written as a paper.  The two activities are not sequential; they are dual.

% ─────────────────────────────────────────────────────────────────────────────
\section{Categorical Treatment: PEI as a Functor}
\label{sec:category}

The isomorphism $\Phi$ has a natural categorical interpretation that deepens
the connection between the document and module worlds.

\subsection{Two Categories}

\begin{definition}[Category of Sovereign Documents $\mathbf{Doc}$]
  $\mathbf{Doc}$ has sovereign research documents as objects.  A morphism
  $f : P \to Q$ is a \emph{citation}: $P$ cites $Q$ as a foundation, expressed
  as a reference in $P_C$ (the sovereignty contract) or as an equation in
  $P_M$ derived from $Q_M$.
\end{definition}

\begin{definition}[Category of Sovereign Modules $\mathbf{Mod}$]
  $\mathbf{Mod}$ has sovereign software modules as objects.  A morphism
  $g : E \to F$ is an \emph{import dependency}: $E$ imports a type or function
  from $F$.
\end{definition}

\begin{theorem}[PEI as a Functor]
  The map $\Phi : \mathbf{Doc} \to \mathbf{Mod}$ defined by
  $\Phi(P) = E$ (PEI applied to objects) and
  $\Phi(f : P \to Q) = (g : \Phi(P) \to \Phi(Q))$ (citation maps to import)
  is a covariant functor.
\end{theorem}

\begin{proof}
  We verify the two functor laws.

  \medskip\noindent
  \textbf{Identity preservation.}  The identity morphism $\text{id}_P : P \to P$
  (a self-citation, i.e.\ a document referencing its own prior section) maps to
  $\Phi(\text{id}_P) : E \to E$ — a module that imports from itself (a recursive
  reference).  This is the identity morphism in $\mathbf{Mod}$.

  \medskip\noindent
  \textbf{Composition preservation.}  For citations $f : P \to Q$ and
  $g : Q \to R$, the composite citation $g \circ f : P \to R$ (``$P$ builds on
  $Q$ which builds on $R$'') maps to the composite import
  $\Phi(g) \circ \Phi(f) : \Phi(P) \to \Phi(R)$ (``module $\Phi(P)$ imports
  from $\Phi(Q)$ which imports from $\Phi(R)$'').  This is exactly
  $\Phi(g \circ f)$.  \qed
\end{proof}

\subsection{Natural Transformation: $\Phi$ and $\Phi^{-1}$}

The inverse functor $\Phi^{-1} : \mathbf{Mod} \to \mathbf{Doc}$ is also a
functor (by symmetric argument).  The pair $(\Phi, \Phi^{-1})$ constitutes a
\emph{category equivalence} — a stronger statement than mere isomorphism of
individual objects.  Category equivalence means not only that each document
corresponds to a module, but that the \emph{dependency structure} of the entire
document corpus is preserved under translation to the module dependency graph.

\begin{corollary}[Knowledge Graph Isomorphism]
  The citation graph of a sovereign document corpus is isomorphic, as a directed
  graph, to the import dependency graph of the corresponding module corpus.
  Software dependency analysis tools can therefore be applied directly to
  document citation analysis, and vice versa.
\end{corollary}

\subsection{The Bidirectional Compiler as an Adjunction}

The $\Phi^{-1}$ compiler (module to paper) and the $\Phi$ compiler (paper to
module) form an adjoint pair.  Specifically, $\Phi^{-1}$ is left adjoint to
$\Phi$: for any document $P$ and module $E$,
\[
  \mathrm{Hom}_{\mathbf{Mod}}(\Phi(P), E) \cong \mathrm{Hom}_{\mathbf{Doc}}(P, \Phi^{-1}(E)).
\]
Intuitively: a module morphism from the module-compiled-from-$P$ to $E$ is the
same as a document morphism from $P$ to the paper-compiled-from-$E$.
This adjunction captures the fact that ``implementing a paper'' and ``documenting
a module'' are two descriptions of the same act.

% ─────────────────────────────────────────────────────────────────────────────
\section{AI Implications: Autonomous Knowledge Production}
\label{sec:ai_implications}

The Paper--Engine Isomorphism has transformative implications for AI systems
that generate, consume, and compose both research and code.

\subsection{LLMs as $\Phi^{-1}$ Compilers}

Large language models trained on code and scientific literature are,
in the PEI framework, approximate implementations of $\Phi^{-1}$: given a
module $E$, an LLM can generate a natural-language description of $E$ that
approximates the paper $P = \Phi^{-1}(E)$.  Similarly, an LLM given a paper
$P$ can generate code that approximates the module $E = \Phi(P)$.

\begin{proposition}[LLM Quality Bound]
  The quality of an LLM as a $\Phi^{-1}$ compiler is bounded by its ability
  to preserve the six structural components $(P_A, P_T, P_M, P_S, P_L, P_C)$.
  A model that generates prose but omits $P_C$ (the sovereignty contract) is
  a lossy compiler that violates Irrevocable Attribution
  (Corollary~\ref{cor:4}).
\end{proposition}

This observation immediately suggests an evaluation benchmark for LLMs as
code documentation tools: does the generated documentation preserve all six
structural components of the source module?  Existing documentation benchmarks
focus on $P_A$ and $P_T$; PEI motivates including $P_M$ (equation extraction),
$P_S$ (state semantics), $P_L$ (lifecycle description), and $P_C$ (attribution).

\subsection{Autonomous Research Agents}

An AI agent equipped with the PEI framework can \emph{produce research
autonomously}:

\begin{enumerate}[label=(\arabic*)]
  \item The agent builds a new module $E$ that solves a problem.
  \item It applies $\Phi^{-1}$ to generate a draft paper $P = \Phi^{-1}(E)$.
  \item It verifies the structural completeness of $P$ (all six components
    present).
  \item It submits $P$ to a paper registry with the module as its proof of
    execution.
\end{enumerate}

Step~2 is where current LLMs excel (code-to-documentation generation).
Step~3 is a structural verification that can be automated.
Step~4 requires an on-chain paper registry that accepts submissions with
linked executable modules.

The result is an AI agent that is simultaneously a software engineer and a
research author — not by metaphor, but by the structural identity established
by PEI.

\subsection{Knowledge Compounding in AI Agent Networks}

In a multi-agent system where agents share a PEI-linked paper and module
registry, knowledge compounds at a rate determined by the composition depth
of the module dependency graph:

\begin{proposition}[Knowledge Compounding Rate]
  In a multi-agent system where each agent can produce one new module per time
  step, and modules can be composed from existing modules in the registry, the
  number of distinct knowledge units accessible to any agent grows as
  $\Theta(B^d)$ where $B$ is the mean branching factor of the dependency graph
  and $d$ is the current composition depth.  The corresponding paper corpus
  grows at the same rate by the functor $\Phi^{-1}$.
\end{proposition}

This exponential growth in accessible knowledge is the agent-network analogue
of the academic compounding effect observed in citation networks: a highly
cited paper enables successively larger bodies of derivative work.  Under PEI,
the same compounding is structural rather than social — it is a consequence of
the module dependency graph, not of human reading habits.

\subsection{Preventing Knowledge Orphaning}

A \emph{knowledge orphan} is a module that exists without a corresponding paper
(a discovery without disclosure) or a paper that exists without a corresponding
module (a disclosure without a running proof).  Both are economically and
epistemologically wasteful.

The PEI framework prevents knowledge orphaning by design: any agent operating
under PEI cannot produce a module without simultaneously producing its paper
skeleton, and cannot publish a paper without linking a running module.  The
attribution record in $E_C = P_C$ closes both directions of the orphan risk.

For AI agent networks where hundreds or thousands of modules are produced
autonomously per day, knowledge orphaning prevention is not a paperwork concern
— it is an existential requirement for maintaining a coherent and auditable
body of machine-generated knowledge.

% ─────────────────────────────────────────────────────────────────────────────
\section{Discussion and Limitations}
\label{sec:discussion}

\subsection{What PEI Does Not Claim}

PEI does not claim that all software is academic research or that all academic
research is executable.  It claims that \emph{sovereign modules built from an
explicit intellectual foundation} and \emph{sovereign documents built from the
same foundation} are isomorphic.  Modules without a clear intellectual
foundation (e.g.\ a configuration file, a boilerplate scaffold) do not satisfy
the definition of a sovereign module and are outside PEI's scope.

\subsection{The Prose Gap}

The $\Phi^{-1}$ direction (module to paper) leaves what we call the
\emph{prose gap}: the narrative that explains \emph{why} the structure is what
it is, not just \emph{what} it is.  This gap cannot be closed by syntactic
rewriting alone; it requires the author to articulate the intellectual journey.
PEI fills the structure; the author fills the story.  For AI agents, the prose
gap is the quality gap between the structural paper skeleton generated by
$\Phi^{-1}$ and a paper good enough for human expert review.

\subsection{Dependency Cycle Risk}

If Paper~A cites Paper~B as a foundation, and module $E_A$ depends on $E_B$,
then a dependency cycle $A \to B \to A$ in either domain implies a cycle in
both — which in a software module would cause a compilation error.  PEI
therefore predicts that a well-formed sovereign knowledge corpus is a
directed acyclic graph (DAG) in the dependency structure, matching the
requirement for acyclic inter-module dependencies in the build system.

\subsection{Evaluation of the Functor Fidelity}

The categorical treatment (Section~\ref{sec:category}) assumes that citations
map faithfully to imports.  In practice, a paper may cite a work that inspired
its approach without the module importing any code from the cited work's module.
These \emph{conceptual citations} (versus \emph{structural citations}) represent
a gap in the functor fidelity.  Distinguishing the two citation types — and
ensuring that structural citations are faithfully represented as imports — is a
discipline requirement for fully realising the categorical isomorphism.

% ─────────────────────────────────────────────────────────────────────────────
\section{Conclusion}
\label{sec:conclusion}

The Paper--Engine Isomorphism establishes that a sovereign research document
and a sovereign software module are two serialisations of the same underlying
abstract object.  The isomorphism is structure-preserving and invertible: every
section maps to a component, and every component maps back to a section.  The
categorical treatment elevates the isomorphism from a bijection on individual
objects to a functor on the entire knowledge corpus, preserving citation
structure, enabling dependency analysis tools to cross between domains, and
establishing the bidirectional compiler as an adjoint pair.

The AI implications are the most consequential: LLMs are approximate
implementations of $\Phi^{-1}$; autonomous research agents can use PEI to
produce papers and modules simultaneously; knowledge compounds at exponential
rate in multi-agent PEI networks; and knowledge orphaning — the disconnection
of discovery from disclosure — is structurally prevented.

The deep implication is this: in a world where AI agents produce capabilities
at machine speed, the rate-limiting factor for knowledge production is not
intelligence but \emph{structure}.  PEI provides that structure.

% ─────────────────────────────────────────────────────────────────────────────
\begin{thebibliography}{99}

\bibitem{curry1934}
H.~B.~Curry,
``Functionality in combinatory logic,''
\textit{Proc.\ Natl.\ Acad.\ Sci.\ USA}, vol.~20, no.~11, pp.~584--590, 1934.

\bibitem{howard1980}
W.~A.~Howard,
``The formulae-as-types notion of construction,''
in \textit{To H.B. Curry: Essays on Combinatory Logic, Lambda Calculus and
Formalism}, Academic Press, 1980.

\bibitem{wadler2015}
P.~Wadler,
``Propositions as types,''
\textit{Commun.\ ACM}, vol.~58, no.~12, pp.~75--84, 2015.

\bibitem{knuth1984}
D.~E.~Knuth,
``Literate programming,''
\textit{Comput.\ J.}, vol.~27, no.~2, pp.~97--111, 1984.

\bibitem{awodey2010}
S.~Awodey,
\emph{Category Theory},
2nd ed., Oxford University Press, 2010.

\bibitem{maclane1998}
S.~Mac~Lane,
\emph{Categories for the Working Mathematician},
2nd ed., Springer-Verlag, 1998.

\bibitem{milner1978}
R.~Milner,
``A theory of type polymorphism in programming,''
\textit{J.\ Comput.\ Syst.\ Sci.}, vol.~17, no.~3, pp.~348--375, 1978.

\bibitem{pierce2002}
B.~C.~Pierce,
\emph{Types and Programming Languages},
MIT Press, 2002.

\bibitem{brown2020}
T.~B.~Brown et al.,
``Language models are few-shot learners,''
\textit{Proc.\ NeurIPS}, 2020.

\bibitem{chen2021}
M.~Chen et al.,
``Evaluating large language models trained on code,''
arXiv:2107.03374, 2021.

\bibitem{dfinity2022}
\text{DFINITY Foundation},
``The Internet Computer,''
Technical White Paper, 2022.

\end{thebibliography}

\end{document}
